\section{物权的法律适用}
\begin{enumerate}
  \item 物权。
        \begin{enumerate}
          \item 基本原则。\zd{物之所在地}。
          \item \zd{不动物之所在地，动产自治获得地}。被追夺方获得该物的事实发生在哪里。
                \begin{kk}{例外}
                  \begin{enumerate}
                    \item 船舶物权纠纷 \yjt 船旗国。\zd{所有权、担保物权纠纷}。例外。
                          \begin{enumerate}
                            \item \zd{优先权}（多物权清偿顺序）\yjt \zd{法院地}。
                            \item \zd{光船（没有挂旗）租赁以前 \hz 期间}设立船舶抵押权的 \yjt \zd{原登记国}。
                          \end{enumerate}
                    \item 民用航空器物权 \yjt 登记国。\zd{所有权、担保物权纠纷}。例外。\zd{优先权} \yjt \zd{法院地}。
                  \end{enumerate}
                  \zd{船旗飞登优法院，光租抵押原登记}。
                  \tcblower
                  \begin{enumerate}
                    \item 运输中动产的物权纠纷 \yjt \zd{运输自治目的地}。
                    \item \zd{证券实现密联地，权利质权设立地}。
                  \end{enumerate}
                \end{kk}
        \end{enumerate}
  \item 债权。
        \begin{enumerate}
          \item 合同之债一般\zd{意思自治 \yjt 最密切联系}。
                \begin{kk}{三种合同允许但限制意思自治}
                  \begin{enumerate}
                    \item \zd{在中国境内履行的中外合资}经营企业合同 \yjt 适用中国法律。\zd{中外合资（履行）在境内}。
                    \item \zd{消费者合同} \yjt 商品、服务提供地。\zd{只能}由消费者选 \yjt 不选 \yjt \zd{消费者经居地有没有营业} \yjt 提供地。\zd{限买限选提供地，不选买居看经营。买居经营买居地，买居无营提供地}。
                    \item \zd{劳动合同}。
                          \begin{enumerate}
                            \item 劳动聘用。工作地 \yjt 主营业地。\zd{劳动合同禁选法，工地确定工作地，工地不定营业地}。
                            \item 劳务派遣。工作地 \hz 派出地 \yjt 主营业地。\zd{派遣还可派出地}。
                          \end{enumerate}
                  \end{enumerate}
                \end{kk}
          \item 侵权之债。意思自治 \yjt 共同经居地 \yjt 侵权行为地。\zd{（普通侵权）自治共居行为地}。
                \begin{kk}{例外}
                  \begin{enumerate}
                    \item 船舶碰撞。\zd{碰在领水侵权地，碰在公海法院地。碰撞限额要注意，只能适用法院地}。
                    \item 民用航空器。\zd{民航领侵公法院}。领土、侵权地、公海、法院地。
                    \item 产品责任。\zd{限弱限选侵权地，不选弱居看经营，弱居经营弱居地，弱居无营侵权地}。
                    \item 人格权。被侵权人及权利人经常居所地。
                    \item 知识产权。允许双方合意选但限制选（见下）\yjt 被请求保护地。
                  \end{enumerate}
                \end{kk}
          \item 不当得利、无因管理之债。\zd{自治共居发生地}。
        \end{enumerate}
  \item 商事关系。
        \begin{enumerate}
          \item 票据。\zd{票据行为行为地，支票出票可协议。追索期限出票地，其他统统付款地}。
          \item 海事和民用航空。共同海损理算适用理算地法。
        \end{enumerate}
  \item 知识产权。\zd{被保护请求地法}。\al{A 与 B 公司就货物在中国境内的商标权产生争议，诉至中国某法院}。
        \begin{enumerate}
          \item \zd{内容归属保护地，转让许可是协议}。
          \item 允许合意，但只能在侵权行为发生后选法院地法；不选，（分别）适用被请求保护地法。
                \al{韩国甲在中韩注册了商标，乙擅自仿冒并销售至中韩两国}。
        \end{enumerate}
\end{enumerate}
